\section{Literature Review}
\subsection{Existing Wireless Foundation Models}

\cite{liu2025wifo} built up on similar ideas  to \cite{alikhani2024large} in developing developing wireless channel models. However, instead of just considering an instance of time, they extend the input channels to 3D ($Time \times Freq \times Antenna$). This can the foundation model to capture temporal changes to different sub carriers frequencies and across different spatial distribution of antennas. Therefore, instead of using 2D patches, 3D patches from which masked tokens are created. Also, instead of an encoder-only architecture, a Masked-AutoEncoder \cite{he2022masked}, which has a transformer decoder in addition to the encoder, to add more capacity in the self-supervised masked reconstruction pretraining. The reconstruction tasks are now either time-masked, frequency masked or a combination of both.
\cite{yang2025wirelessgpt} proposed WirelessGPT with the goal of buidling a foundation model that extends beyond physical layer communication tasks but also to sensing tasks. The sensing downstream tasks are human activity recognition and wireless environment reconstruction. The human activity recognition involves using CSI to classify user activities such as running, walking etc. The environment reconstruction involves producing 3D point cloud representations of the environment based on the model's understanding of shadowing, scattering and reflection of rays in the environment inferred from the CSIs. 

\cite{yu2024channelgpt} introduced \textit{ChannelGPT}, a digital twin framework designed to integrate multimodal channel models into emerging 6G networks. The framework emphasized the adoption of feedback-loops both from changes to the environment and network performance to guide the operations of the channel models. Also,     the framework mentions the need of utilizing multimodal data sources, such as multi-view images of the communication environment together with channel state information (CSI), to achieve a richer and more context-aware representation of the wireless channel. The authors argue that incorporating multimodal data enables a deeper understanding of the relationship between the physical environment and the channel characteristics, thereby enhancing network performance.In contrast, other models assume that this environment and channel relationship can be implicitly learned from the variations present in the CSI data alone, without explicitly relying on multimodal inputs.